% !PRINTED 401
\weberpage{401}

3) Es seien die gegebenen Linien:
\[
0,\ 1,\ 2,\ 3,\ 4,\ 5,\ 5',
\]
und davon 0, 1, 2, 3, 4 reell, 5 und 5' conjugirt imagin\"ar. Die
Betrachtung der drei Complexe:
\[
\begin{aligned}
(0)\quad & 1\,[1\,0],\ 2\,[2\,0],\ 3\,[3\,0],\ 4\,[4\,0],\ 5\,[5\,0],\ 5'\,[5'\,0] \\
(5)\quad & 0\,[0\,5],\ 1\,[1\,5],\ 2\,[2\,5],\ 3\,[3\,5],\ 4\,[4\,5],\ 5'\,[5'\,5] \\
(5')\quad & 0\,[0\,5'],\ 1\,[1\,5'],\ 2\,[2\,5'],\ 3\,[3\,5'],\ 4\,[4\,5'],\ 5\ [5\,5'],
\end{aligned}
\]
von denen der erste reell, die beiden anderen conjugirt imagin\"ar
sind, und der durch Vertauschung von 0, 1, 2, 3, 4 aus (0) abgeleiteten
 giebt, genau wie oben, folgendes Resultat:

\begin{center}
\gesperrt{Reelle Doppeltangenten:}
\end{center}
\[
\begin{aligned}
&0,\quad 1,\quad 2,\quad 3,\quad 4,\quad [0\,1],\ [0\,2],\ [0\,3] \\
&[0\,4]\ [1\,2],\ [1\,3]\ [1\,4],\ [2\,3],\ [2\,4],\ [3\,4],\ [5\,5'].
\end{aligned}
\]

\begin{center}
\gesperrt{Conjugirte Paare:}
\end{center}
\[
\begin{aligned}
&5,\ [0\,5],\ [1\,5],\ [2\,5],\ [3\,5],\ [4\,5] \\
&5',\ [0\,5'],\ [1\,5'],\ [2\,5'];\ [3\,5'],\ [4\,5'], % [sic] semicolon printed where a comma is expected
\end{aligned}
\]
und wir finden ein syzygetisches Complextripel:
\[
\begin{aligned}
&0,\quad 1,\quad [0\,2][1\,2],\ [0\,3][1\,3],\ [0\,4][1\,4\,],\ [0\,5][1\,5\,],\ [0\,5'][1\,5'] \\
&2,\quad 3,\quad [0\,2][0\,3],\ [1\,2][1\,3],\ [4\,2][4\,3\,],\ [2\,5][3\,5\,],\ [2\,5'][3\,5'] \\
&[0\,1],\ [2\,3],\ [0\,2][1\,3],\ [0\,3][1\,2],\quad 4\ \ [5\,5'],\quad 5\ \ [4\,5'],\quad 5'\ \ [4\,5],
\end{aligned}
\]
von denen jeder Complex noch vier reelle Paare und ein imagin
\"ares Doppelpaar enth\"alt. Solcher Tripel lassen sich aber noch
mehrere bilden, da man 0, 1, 2, 3, 4 beliebig vertauschen darf.
Dies ist der dritte Fall von \S.\,104.

4) Dass endlich, wenn wir alle Elemente des gegebenen
Siebener-Systemes reell voraussetzen, auch alle \"ubrigen Doppeltangenten
 reell ausfallen, ist eine unmittelbare Folge der rationalen
 Darstellung (\S.\,100).

Hiermit ist gezeigt, dass die vier F\"alle des vorigen Paragraphen
 wirklich alle vorkommen k\"onnen. Ob die hier besprochene
 Erzeugungsweise die einzig m\"ogliche ist, mit anderen
Worten, ob bei jeder reellen Curve vierter Ordnung ein reelles
\gesperrt{Aronhold'sches} System existirt, diese Frage m\"ussen wir unentschieden
 lassen.
